\documentclass{texlr}

\title{Texlr kitchen sink: every class feature on paper}
\author{Texlr test agent}
\date{30 August 2026}
\project{Texlr}
\status{Visual regression fixture}
\repository{https://github.com/drdreo/texlr}
\revision{kitchen-sink}
\runningtitle{Texlr kitchen sink}

\begin{document}
\maketitle

\section{Purpose}

This document exercises every feature of the \texttt{texlr} class in one
build so a rendered copy can be compared against a previous screenshot. Each
section states what to look for. The title block above should show all six
metadata rows in monochrome sans-serif type.

\section{Typography and inline markup}

Body text is serif with a \(1.5\)\,em paragraph indent and no inter-paragraph
gap. Inline markup: \emph{emphasis}, \textbf{bold}, \texttt{inline-code},
a breakable path like \path{internal/texlr/templates/texlr.cls}, and a URL
such as \url{https://github.com/drdreo/texlr/blob/main/README.md} rendered in
sans-serif without a colored border.

Escaped specials must all print literally: \& \% \$ \# \_ \{ \}
\textasciitilde{} \textasciicircum{} \textbackslash{} and a
footnote\footnote{Footnotes use the article default layout.} completes the set.

Inline math \(e^{i\pi} + 1 = 0\) and a displayed equation:
\begin{equation}
  \operatorname{coverage}(D) = \frac{\lvert \text{features shown} \rvert}
                                    {\lvert \text{features total} \rvert}
  \label{eq:coverage}
\end{equation}
Equation~\ref{eq:coverage} also verifies cross-references.

\section{Lists}

\subsection{Itemize and enumerate}

\begin{itemize}
  \item First unordered entry with tight vertical spacing.
  \item Second entry.
  \begin{itemize}
    \item Nested entry.
  \end{itemize}
\end{itemize}

\begin{itemize}[label=$\square$]
  \item Intentional unchecked task marker.
\end{itemize}

\begin{enumerate}
  \item First ordered step.
  \item Second ordered step.
\end{enumerate}

\subsection{Description}

\begin{description}
  \item[callout] A bordered panel with an accent-colored title.
  \item[decision] A callout preloaded with a ``Decision:'' title.
\end{description}

\section{Callouts and decisions}

\begin{callout}[Verify first]
The callout title above must render in the blue accent while this body text
stays black on the light panel background with a visible border.
\end{callout}

\begin{callout}
A callout without an optional argument falls back to the ``Note'' title.
\end{callout}

\decision{Adopt the kitchen-sink fixture}{Keeping one document that shows
every feature makes visual regressions reviewable from a single build.}

\section{Table}

\begin{table}[ht]
  \centering
  \caption{Feature coverage summary.}
  \label{tab:coverage}
  \begin{tabularx}{\linewidth}{@{}l X l@{}}
    \toprule
    Feature & What to check & State \\
    \midrule
    Title block & Six metadata rows, monochrome & Shown \\
    Callout & Accent title, panel background & Shown \\
    Diagrams & Graphviz and Mermaid figures & Shown \\
    Code & Numbered listing on gray background & Shown \\
    \bottomrule
  \end{tabularx}
\end{table}

Table~\ref{tab:coverage} uses booktabs rules and a stretching X column.

\section{Code listing}

\begin{lstlisting}[language=bash,caption={Build command under test}]
texlr build kitchen-sink.tex \
  --pdf artifacts/kitchen-sink.pdf \
  --source artifacts/source \
  --log artifacts/kitchen-sink.log --force
\end{lstlisting}

\section{Figures}

\subsection{Graphviz}

\begin{figure}[ht]
  \centering
  \graphviz[width=0.85\linewidth]{diagrams/pipeline.dot}
  \caption{Texlr build pipeline rendered from a DOT source.}
  \label{fig:pipeline}
\end{figure}

Figure~\ref{fig:pipeline} should use the restrained blue-gray palette.

\subsection{Mermaid}

\begin{figure}[ht]
  \centering
  \mermaid[width=0.9\linewidth]{diagrams/states.mmd}
  \caption{Build states rendered from a Mermaid source.}
  \label{fig:states}
\end{figure}

\subsection{Ordinary image}

\begin{figure}[ht]
  \centering
  \includegraphics[width=0.5\linewidth]{images/sample.png}
  \caption{A PNG included with ordinary \texttt{\textbackslash includegraphics}.}
  \label{fig:png}
\end{figure}

\section{Page furniture}

From page two onward the header must show the short running title
``Texlr kitchen sink'' on the left and the page number on the right, with no
rule. Page one uses the plain style with a centered footer number.

\end{document}
